\section{Revision: lacunary variation of the weighted averages}
The source supplied on 2 November 2026 adds Corollary 4.2, labelled
\texttt{c:long-variation}. The averages remain the weighted averages
$B_N^\gamma$ and $B_N$ of Section 4. No variational assertion for the
count-normalized unweighted averages is claimed. The admissible $c_*$
is selected before the variation exponent and the lacunarity ratio.

\begin{definition}[Infinite variation seminorm]\label{def:variation}
\lean{PSPrimes.variationChain, PSPrimes.variationChain_eq_ofReal, PSPrimes.rVariation}\leanok
For $1\leq r<\infty$ and $a_j\in\C$, set
\[
 V^r(a)=\sup_{J\geq1,\;j_0<\cdots<j_J}
 \left(\sum_{h=1}^J|a_{j_h}-a_{j_{h-1}}|^r\right)^{1/r}.
\]
The Lean definition takes this supremum in $[0,\infty]$, so infinite
variation is not assigned an arbitrary real value. Sequence indices start
at zero in Lean; shifting them by one gives exactly the manuscript's
indexing. The extra $J=0$ term is zero and does not change the supremum.
\end{definition}
\Status{checked definition and finite-chain correspondence with the real manuscript formula}
\ifDetailed
For $j:\{0,\ldots,J\}\to\N$, \texttt{variationChain} is the direct
extended-real expression
\[
 \left(\sum_{h=0}^{J-1}
  \|a_{j(h+1)}-a_{j(h)}\|_e^r\right)^{1/r}.
\]
\texttt{variationChain\_eq\_ofReal} proves equality with the paper's real finite-chain formula. The supremum ranges over $J$ and all such functions, with value zero when
$j$ is not strictly increasing. This is a countable supremum. Hence
measurability follows from measurability of each finite expression;
almost-everywhere measurability is obtained by choosing measurable versions
on one common conull set. This countability is needed for a genuine
$L^1$ statement about the supremum, rather than estimates for each fixed
chain separately. Extended nonnegative integrals avoid the convention
that a nonintegrable real-valued function has integral zero.
\fi

\begin{auxlemma}[Summable perturbations and lacunary scales]\label{l:variation-perturbation}
\uses{def:variation}
\lean{PSPrimes.rVariation_add_le, PSPrimes.rVariation_le_two_tsum, PSPrimes.lintegral_rVariation_le_two_tsum, PSPrimes.lacunary_rpow_tsum_ennreal_le}\leanok
For $r\geq1$, complex sequences satisfy
\[
 V^r(a+b)\leq V^r(a)+V^r(b),\qquad
 V^r(a)\leq2\sum_{j\geq0}|a_j|.
\]
If $N_j\in\N_{>0}$ and $N_{j+1}\geq\lambda N_j$ with
$\lambda>1$, then for every $\eta>0$,
\[
 \sum_{j\geq0}N_j^{-\eta}
 \leq\sum_{j\geq0}\lambda^{-j\eta}
 =\frac1{1-\lambda^{-\eta}}.
\]
Consequently, for almost-everywhere measurable $E_j:X\to\C$,
\[
 \int_X V^r(E_j(x))\,d\nu(x)
 \leq2\sum_{j\geq0}\int_X |E_j(x)|\,d\nu(x).
\]
\end{auxlemma}
\Status{checked perturbation, measurability and infinite geometric-sum proofs}
\begin{proof}[Perturbation and geometric summation]
\proves{l:variation-perturbation}\uses{def:variation}\leanok
\ifDetailed
Minkowski's inequality for each finite chain proves the triangle
inequality, and taking the countable supremum preserves it. For the
second inequality, the finite $\ell^r$ norm is at most the $\ell^1$ norm
when $r\geq1$. The triangle inequality bounds each difference by its two
endpoint norms. Strict increase makes each left-endpoint map and each
right-endpoint map injective, so their sums are separately at most
$\sum_j|a_j|$. This explains the constant $2$ without imposing adjacent
chain indices. Tonelli's theorem then moves the nonnegative infinite sum
through the integral.

For scales, induction gives $N_j\geq\lambda^jN_0\geq\lambda^j$.
The function $t\mapsto t^{-\eta}$ is decreasing for $t>0$, and
$0<\lambda^{-\eta}<1$. Comparison with the geometric series proves both
summability and the displayed bound. The constant does not depend on
$N_0$ or on any later member of the sequence. This is a bound on the
infinite sum, not a finite-range numerical check.
\else
Apply finite-dimensional Minkowski, bound each chain endpoint sum by the
full nonnegative sum, and use Tonelli. Induction gives $N_j\geq\lambda^j$;
comparison with the geometric series gives the uniform bound.
\fi
\end{proof}

\begin{inputthm}[Ordinary-prime lacunary variation]\label{input:MWV}
\uses{def:averages,def:variation}
\lean{PSPrimes.OrdinaryPrimeVariationInput}
For each fixed distinct-degree integer polynomial family, $r>2$ and
$\lambda>1$, there is a finite constant $C_{\mathbf P,r,\lambda}$ such
that, uniformly over probability systems, essentially $1$-bounded inputs
and positive $\lambda$-lacunary sequences,
\[
 \|V^r(B_{N_j}:j\geq0)\|_1\leq C_{\mathbf P,r,\lambda}.
\]
\end{inputthm}
\Status{external MW-V, not a project proof}
Mirek--Wan--Wright, arXiv:2505.15549v2, Theorem 1.1(iv), equation (1.5),
provides the published estimate. Its source-to-interface translation
MW-V-B is recorded separately in the shared catalogue. Ordinary pointwise
convergence alone does not supply this bound.

\begin{theorem}[Lacunary variation; manuscript Corollary 4.2]\label{c:long-variation}
\uses{def:objects,def:averages,def:variation,l:variation-perturbation,input:MWV,t:shao-sun,l:orbit-transference}
\lean{PSPrimes.lacunary_variation_of_external_inputs, PSPrimes.orbit_variation_bound_eLpNorm_of_oneBounded}\leanok
For the fixed family $\mathbf P$, there is $c_*>1$ such that, for every
$1<c<c_*$, $r>2$ and $\lambda>1$, every probability-preserving invertible
system and $f_i\in L^\infty(X)$ satisfy
\[
 \left\|V^r(B_{N_j}^\gamma:j\geq0)\right\|_1
 \leq C_{c,\mathbf P,r,\lambda}\prod_{i=1}^\ell\|f_i\|_\infty
\]
for every sequence of positive integers with
$N_{j+1}\geq\lambda N_j$. Here $\gamma=1/c$, and the constant is chosen
before the sequence, system and functions.
\end{theorem}
\Status{checked conditional proof from SS, TV and MW-V}
\begin{proof}[Transfer argument]
\proves{c:long-variation}\leanok
\uses{l:variation-perturbation,input:MWV,t:shao-sun,l:orbit-transference}
For essentially $1$-bounded inputs put $E_j=B_{N_j}^\gamma-B_{N_j}$.
The already checked all-scale estimate
\texttt{ps\_orbit\_discrepancy\_L1\_gamma} supplies $D>0$ and
$\eta>0$ with $\|E_j\|_1\leq DN_j^{-\eta}$. Hence
\[
 \|V^r(E_j)\|_1\leq
 \frac{2D}{1-\lambda^{-\eta}}.
\]
Apply MW-V to $B_{N_j}$ and use the variation triangle inequality.
For general inputs divide each nonzero input by its essential norm and
use multilinearity. If one essential norm is zero, that input vanishes
almost everywhere; measure preservation and countability give vanishing
of every sampled product on a common conull set. Thus the zero-factor
case also satisfies the stated product bound.
\end{proof}
\ifDetailed
The choice of $c_*$ and the saving $\eta$ comes from SS and TV, before
$r$ and $\lambda$. For fixed $c,r,\lambda$, one may take the unit-input
constant
\[
 C_{\mathbf P,r,\lambda}^{\rm MWV}
 +\frac{2D_{c,\mathbf P}}{1-\lambda^{-\eta}}.
\]
The ordinary variational estimate is the only new substantial external
input. No PNT, weighted mass limit, a.e. pointwise conclusion,
subsequence interpolation, removal of higher prime powers or Abel
unweighting is needed for this weighted corollary. It does not give a
full-sequence variation estimate, an $r=2$ estimate, or a finite-$L^p$
extension of the sparse-prime theorem. The new endpoint Remark 1.3 makes
such an extension an expectation only.
\fi
